% !TEX program = xelatex
\documentclass[11pt,a4paper]{ctexart}
\usepackage[a4paper,margin=2.2cm]{geometry}
\usepackage{amsmath,amssymb}
\usepackage{booktabs}
\usepackage{enumitem}
\usepackage{hyperref}
\usepackage{listings}
\usepackage{xcolor}
\usepackage{longtable}
\usepackage{array}

\hypersetup{colorlinks=true,linkcolor=blue,urlcolor=blue,citecolor=blue}

\title{显式静电 PIC 模块通用化：实现路径与可复用架构设计\\\large（优先：量纲体系、模块字典规范、MCC 碰撞/反应；面向隐式 PIC 与电磁扩展）}
\author{陈冲（Chong Chen）}
\date{2026-01-23}

\lstset{basicstyle=\ttfamily\small,breaklines=true,frame=single,rulecolor=\color{black!20}}

\begin{document}
\maketitle

\begin{abstract}
本文给出一个可执行的“从现有 PaperA 显式静电 PIC 流程（deposit\,$\rightarrow$\,Poisson\,$\rightarrow$\,E\,$\rightarrow$\,gather\,$\rightarrow$\,push\,$\rightarrow$\,output）到通用可用模块”的实现路径。
考虑到后续要扩展隐式 PIC、电磁（EM）与更多物理模块，路线以“先抽出可复用基础设施”为原则：\textbf{量纲/归一化（Units）、模块字典规范与校验（Schema/Validator）、碰撞/反应 MCC 框架（Cross section + Reaction DSL + Operator 接口）}优先落地；随后再补齐边界/源项/诊断/算例库，最终形成可回归、可复现、可扩展的模块体系。
\end{abstract}

\tableofcontents

\section{总体目标与设计原则}
\subsection{目标（验收口径）}
面向 “IEC 分析 + CPC/经典 benchmark + PaperA 算例”，模块达到：
\begin{itemize}[leftmargin=2em]
\item \textbf{输入可表达真实工况：}几何/网格由 \texttt{geometry.dict} 与 \texttt{mesh.dict} 描述；显式 PIC 模块字典可用 SI 输入（可选内部归一化）。
\item \textbf{物理闭环可配置：}物种、源项、壁面、碰撞/反应、外电路（至少浮动电极/电流平衡）可通过字典组合。
\item \textbf{诊断与回归完备：}能量、Gauss 残差、壁面电流、反应率、能谱等作为 ``golden outputs'' 可回归。
\item \textbf{架构可扩展：}同一套 Units/字典/碰撞/诊断框架可复用于隐式 PIC 与 EM 模块。
\end{itemize}

\subsection{核心原则}
\begin{enumerate}[leftmargin=2em]
\item \textbf{先基础设施、后堆物理：}先把 Units + Dict Schema + MCC 框架做扎实，再做更多边界/源项/算例。
\item \textbf{接口先行：}Deposition/Gather、Pusher、FieldSolver、CollisionOperator、BoundaryModel、SourceModel 都以接口形式落地，显式/隐式/EM 仅替换实现。
\item \textbf{算例驱动：}每一步都有最小验收算例（benchmark）与固定随机种子，防止“跑得动但变味”。
\end{enumerate}

\section{建议的可复用架构分层（面向后续隐式/电磁扩展）}
\subsection{分层与职责}
建议将 PIC 相关复用能力沉淀为 ``PIC Core''（不依赖具体显式/隐式算法），并由不同 Module 组合：
\begin{itemize}[leftmargin=2em]
\item \textbf{Core/IO：}字典解析、Schema 校验、单位解析与打印、日志。
\item \textbf{Core/Units：}SI \,$\leftrightarrow$\, CodeUnits 的转换与参考尺度管理。
\item \textbf{PICCore/Particles：}物种、宏粒子权重、分布采样、RNG、重启序列化。
\item \textbf{PICCore/Operators：}Deposit/Gather、Pusher、CollisionOperator、SourceOperator、BoundaryOperator。
\item \textbf{FieldCore/Solvers：}Poisson（ES）、Maxwell（EM）等 ``FieldSolver'' 接口与实现；线性求解器统一接口。
\item \textbf{Diagnostics：}全局守恒、残差、壁面通量、能谱、反应统计；统一输出到现有 OutputManager。
\end{itemize}

\subsection{关键接口（建议最小集合）}
\begin{lstlisting}
// 统一单位/尺度
struct UnitsContext { /* SI<->code, reference scales, printing */ };

// 模块运行时上下文（可被不同模块共享）
struct SimContext {
  const Mesh& mesh;
  UnitsContext units;
  RNG rng;
  Diagnostics diag;
  // ... output, timers, etc.
};

// 场求解器（ES/EM/Implicit都能复用）
class IFieldSolver {
public:
  virtual void assemble(const SimContext&, const FieldRHS& rhs) = 0;
  virtual void solve(SimContext&, FieldState& state) = 0;
  virtual ~IFieldSolver() = default;
};

class IDepositOperator { /* particles -> grid rho */ };
class IGatherOperator  { /* grid E/phi -> particles */ };
class IPusher          { /* advance x,v */ };
class ICollisionOperator { /* MCC events */ };
class ISourceOperator  { /* inject/ionize/recombine */ };
class IBoundaryOperator{ /* wall interaction + current/SEE */ };
\end{lstlisting}

\section{实现路径（一步一步）：优先 Units、字典规范、MCC}
下面按 ``Step 0--12'' 给出推荐顺序。每一步都包含：\textbf{目的证明（Why）、交付物（Deliverables）、验收算例（Acceptance）}。

\subsection*{Step 0：冻结基线与回归骨架（1 次性）}
\textbf{Why：}避免后续重构导致 ``能跑但结果漂移''。

\textbf{Deliverables：}
\begin{itemize}[leftmargin=2em]
\item 建立 \texttt{cases/} 目录：\texttt{paperA/}, \texttt{bench/}, \texttt{iec/} 三类。
\item 统一 ``run script''（本地/CI 均可），固定 RNG seed。
\item 现有流程输出最小诊断：总电荷、Gauss 残差、场能/动能。
\end{itemize}
\textbf{Acceptance：}PaperA 至少 1 个算例跑通并生成基线曲线文件（作为 golden）。

\subsection*{Step 1：Units/Normalization 工具层（优先）}
\textbf{Why：}真实工况输入必须有量纲；内部可选无量纲以保持数值健康。

\textbf{Deliverables：}
\begin{enumerate}[leftmargin=2em]
\item \textbf{UnitsContext}：支持 SI 与 CodeUnits 双模式；记录 $L_0,T_0,\Phi_0,n_0,v_0,E_0,\rho_0$ 等参考量。
\item 统一 ``单位解析''：字典字段支持 ``\texttt{value + unit}''（例如 \texttt{50 mTorr}、\texttt{3 cm}、\texttt{500 V}）。
\item 统一输出：日志打印完整尺度表；输出文件中同时写 SI 与 code 值（可选）。
\end{enumerate}
\textbf{Acceptance：}在同一算例中，用 SI 输入与用旧 code 输入得到一致的无量纲结果（误差在统计噪声内）。

\subsection*{Step 2：模块字典规范（Schema）与严格校验（优先）}
\textbf{Why：}通用模块必须有清晰的配置契约；并且要为后续隐式/EM 扩展预留结构。

\textbf{Deliverables：}
\begin{enumerate}[leftmargin=2em]
\item 定义 \texttt{ExplicitPIC.dict} 规范（见下节建议结构），并提供：
  \begin{itemize}[leftmargin=2em]
  \item 必填/可选字段、默认值、取值范围（例如形函数类型、子循环次数、MCC 开关等）。
  \item ``严格白名单''：未知字段报错（或至少警告）。
  \end{itemize}
\item 实现 Schema/Validator：对字段类型（scalar/vector/block）、单位、范围进行校验。
\item 版本化：\texttt{schemaVersion} 字段，允许平滑升级。
\end{enumerate}
\textbf{Acceptance：}对一组故意写错的 dict（缺字段/错单位/错类型）能给出清晰报错；对旧 PaperA dict 能通过兼容层迁移（打印 deprecations）。

\subsection*{Step 3：物种与常数数据库（为 MCC/隐式/EM 打底）}
\textbf{Why：}MCC、迁移率、推导尺度都需要标准化的物种参数与常数。

\textbf{Deliverables：}
\begin{itemize}[leftmargin=2em]
\item \texttt{Species}：$m,q$、宏粒子权重、初始分布、温度/漂移、是否为电子/离子/中性。
\item 常数表：$\epsilon_0,\mu_0,k_B,e$ 等。
\item 统一 ``气体/背景'' 描述：气体种类、压力、温度、密度计算（理想气体）。
\end{itemize}
\textbf{Acceptance：}给定 $p,T$ 自动算出 $n_{gas}$ 并打印；与手算一致。

\subsection*{Step 4：MCC 框架骨架（优先）}
\textbf{Why：}IEC 与放电类 benchmark 若无碰撞/反应，很难对标；框架需尽早定型以复用到隐式 PIC。

\textbf{Deliverables：}
\begin{enumerate}[leftmargin=2em]
\item CrossSectionDB：支持 ``能量 eV \,$\rightarrow$\, 截面'' 曲线表 + 插值（线性/对数）。
\item Reaction DSL：用字典定义反应通道（elastic/excitation/ionization/CX），并绑定产物与能量分配模型。
\item ICollisionOperator：实现 null-collision/majorant 方案；对每个粒子给出 ``是否发生事件'' 与 ``事件类型''。
\item 统计器：反应计数、反应率、能量损失；输出到 Diagnostics。
\end{enumerate}
\textbf{Acceptance：}在均匀气体箱算例中，电子的平均碰撞频率与理论 $\nu=n\sigma v$ 一致（统计误差内）。

\subsection*{Step 5：最小反应集落地（e--neutral 弹性 + 电离）}
\textbf{Why：}这是 ``能跑放电/IEC 的最小集合''。

\textbf{Deliverables：}
\begin{itemize}[leftmargin=2em]
\item e--neutral elastic：散射角模型（各向同性/各向异性可选），能量更新。
\item e--neutral ionization：生成二次电子 + 离子；能量分配（阈值 + 随机分配）。
\item 粒子生成与注入统一走 SourceOperator（MCC 产生也算一种 source）。
\end{itemize}
\textbf{Acceptance：}Townsend-like 体电离算例（简化）：给定均匀电场，电离率随场强增大趋势正确；反应统计可复现。

\subsection*{Step 6：模块字典扩展到 ``Sources + Boundaries + Circuits''（与 MCC 同期）}
\textbf{Why：}MCC 产生的粒子、边界发射、外部束注入都需要统一配置结构。

\textbf{Deliverables：}
\begin{itemize}[leftmargin=2em]
\item Sources：\texttt{volumeInject}, \texttt{surfaceEmit}, \texttt{beamInject}, \texttt{mccIonization}（内部）。
\item Boundaries：粒子-壁面模型（吸收/反射/再发射/SEE）；场边界类型（Dirichlet/Neumann/Robin/Floating）。
\item Circuits：最小 ``CurrentBalance''（浮动电极）与 ``VoltageDrive''（给定波形）。
\end{itemize}
\textbf{Acceptance：}1D sheath/diode：输出 $I(t)$ 与 IEDF；浮动电极电位收敛到合理值（电流平衡）。

\subsection*{Step 7：诊断体系升级（与回归指标固化）}
\textbf{Why：}可分析性来自诊断，而不仅是云图。

\textbf{Deliverables：}
\begin{itemize}[leftmargin=2em]
\item 全局：总电荷、Gauss 残差（L1/L2/Linf）、能量（场能/动能/总能）与漂移率。
\item 边界：按边界/物种的电流与能量通量；IEDF（分箱统计）。
\item 反应：每通道反应率/能量损失；空间分布（可选）。
\item 输出规范：统一文件命名与 units 元数据（便于后处理）。
\end{itemize}
\textbf{Acceptance：}plasma oscillation：频率与 $\omega_{pe}$ 一致；能量漂移在可控范围。

\subsection*{Step 8：AMR/粗细界面鲁棒性（paperA/paperB 风格问题前置防线）}
\textbf{Why：}若后续要做隐式 PIC/更复杂算例，AMR 伪影（自力、界面噪声）会反复出现。

\textbf{Deliverables：}
\begin{itemize}[leftmargin=2em]
\item composite deposition / correction band 的落地（与你已讨论方向一致）。
\item 在 coarse--fine 界面输出专项诊断：局部 Gauss 残差、局部能量与粒子统计。
\end{itemize}
\textbf{Acceptance：}自力检查算例：单粒子/对称分布在粗细界面附近无系统性加速（统计期望接近 0）。

\subsection*{Step 9：重启/检查点（工程闭环）}
\textbf{Why：}长时间积分与参数扫必备；也是回归与调试的基础。

\textbf{Deliverables：}
\begin{itemize}[leftmargin=2em]
\item 序列化：粒子（x,v,w,label）、场（phi,E,rho）、RNG state、时间步信息。
\item 重启一致性：重启后统计量连续，且与不重启版本一致（误差在噪声内）。
\end{itemize}
\textbf{Acceptance：}任意步写 checkpoint，重启继续跑，关键诊断曲线无跳变。

\subsection*{Step 10：IEC 专用能力补齐（面向 \texttt{iec\_pic} 对标）}
\textbf{Why：}IEC 需要更完整的电极/网格/注入/碰撞组合与特定诊断。

\textbf{Deliverables：}
\begin{itemize}[leftmargin=2em]
\item IEC 几何与边界模板（球/圆柱网格，栅极/阴极/阳极配置）。
\item 离子注入与中性背景、CX 通道（若 IEC 工况需要）。
\item IEC 输出：通过栅极电流、离子能谱、陷获率/损失率、空间电位结构。
\end{itemize}
\textbf{Acceptance：}复现 \texttt{iec\_pic} 中至少 1 个代表性工况的主要曲线/统计指标。

\subsection*{Step 11：为隐式 PIC 预留 ``隐式接口替换点''}
\textbf{Why：}隐式 PIC 不是在显式上堆补丁，而是替换 ``FieldSolver + Particle Advance + Coupling''。

\textbf{Deliverables：}
\begin{itemize}[leftmargin=2em]
\item 将 ``显式时间推进'' 写成可替换的 \texttt{IAdvanceScheme}（ExplicitLeapfrog / ImplicitECSIM 等）。
\item 将 ``沉积/收集'' 抽象为可复用算子（隐式也要 matched deposit/gather）。
\item 线性系统接口统一（隐式需要更多块矩阵/耦合项，但求解器接口不变）。
\end{itemize}
\textbf{Acceptance：}在不改 cases 的前提下，显式推进仍能跑；新接口对调用方透明。

\subsection*{Step 12：为 EM（Maxwell）扩展预留 ``场变量与求解器族''}
\textbf{Why：}EM 需要 $\mathbf{E},\mathbf{B}$ 与 Maxwell 求解器，但 PICCore 的粒子/碰撞/源/边界很多可复用。

\textbf{Deliverables：}
\begin{itemize}[leftmargin=2em]
\item FieldState 抽象：ES（phi,E）与 EM（E,B, potentials）统一容器。
\item IFieldSolver 支持 ``ESPoissonSolver'' 与 ``EMMaxwellSolver'' 并存。
\item 边界模型拆分：场边界（电磁边界/导体）与粒子边界（吸收/SEE）组合。
\end{itemize}
\textbf{Acceptance：}最小 EM 测试（例如真空中波传播/能量守恒）可作为未来里程碑，但接口现在就要能承载。

\section{字典规范建议（面向扩展：ES/Implicit/EM 共用骨架）}
推荐将 PIC 模块字典分为 ``通用块 + 具体算法块''：
\begin{lstlisting}
PIC
{
  schemaVersion 1;
  units { mode SI; normalization problem; L0 3 cm; Phi0 500 V; n0 1e15 1/m^3; }

  species
  {
    e { charge -1 e; mass 1 me; weight 1e8; }
    Arp { charge +1 e; mass 40 amu; weight 1e8; }
  }

  background { gas Ar; pressure 50 mTorr; temperature 300 K; }

  numerics
  {
    shapeFunction CIC;     // NGP|CIC|TSC
    subcycling { e 10; Arp 1; }
    filters { rho none; E none; }
    diagnostics { energy true; gauss true; wallCurrent true; IEDF true; }
  }

  boundaries
  {
    cathode { field { type Dirichlet; phi 0 V; }
              particles { type Absorb; SEE { enabled true; model simple; yield 0.1; } } }
    anode   { field { type Dirichlet; phi 500 V; }
              particles { type Absorb; } }
    wall    { field { type Floating; circuit { type CurrentBalance; } }
              particles { type DiffuseReflect; T 300 K; } }
  }

  sources
  {
    beam1 { type BeamInject; species Arp; rate 1e18 1/s; energy 100 eV; }
  }

  collisions
  {
    enabled true;
    scheme nullCollision;
    reactions
    {
      e_Ar_elastic { type Elastic; projectile e; target Ar; xsTable "xs/e-Ar-elastic.dat"; }
      e_Ar_ion     { type Ionization; projectile e; target Ar; threshold 15.76 eV;
                     products (e Arp); xsTable "xs/e-Ar-ion.dat"; }
    }
  }

  algorithm
  {
    family electrostatic;   // electrostatic|implicit_es|electromagnetic
    timeIntegration leapfrog; // leapfrog|boris|ecsim|... (future)
  }
}
\end{lstlisting}

\section{推荐的最短落地顺序（浓缩版）}
如果你希望 ``尽快能做实际问题''，同时又不牺牲后续扩展，建议把工作压成四个短迭代：
\begin{enumerate}[leftmargin=2em]
\item \textbf{Iter-1（基础设施）：}Step 0--2（回归骨架 + Units + Dict Schema）。
\item \textbf{Iter-2（MCC 最小可用）：}Step 3--5（物种/常数 + MCC 框架 + 电离最小集）。
\item \textbf{Iter-3（可分析）：}Step 6--7（Sources/Boundaries/Circuits + 诊断固化）。
\item \textbf{Iter-4（可对标）：}Step 8--10（AMR 鲁棒 + 重启 + IEC 对标算例）。
\end{enumerate}
完成 Iter-4 后，显式 PIC 基本具备 ``通用可用 + 可对标 + 可扩展'' 的工程形态；Step 11--12 则保证你在做隐式/EM 时不会推倒重来。

\section{结语}
你倾向把 ``量纲、字典规范、MCC'' 放前面是正确的：它们决定了整个系统的 ``可表达性、可配置性与可扩展性''。
按本路线推进，可以确保显式 PIC 不是一次性 demo，而是后续隐式 PIC 与电磁模块的可复用基座。

\end{document}
